% GENERATED FILE. Edit canonical lessons, then run npm run generate:books.
% canonical-source-hash: fnv1a64:9443aef900711b7c
% canonical-lessons: BN-W04-chaa-read, BN-W04-jol-read, BN-W03-dudh-read

\chapter{Three Words, No New Pieces}
\label{ch:bn-three-words-no-new-pieces}
\hlchaptermodality{15}

\begin{chapteropening}
\textbf{By the end of this chapter:} \emph{I can read \bn{চা}, \bn{জল} and \bn{দুধ} off the page without learning a single new shape, and say why a reading chapter that adds nothing is still worth its place.}

Milk puts the plain dental and its breathy partner side by side in one three-piece word, then the chapter's three words are read back — every one of them paid for by consonants bought several chapters earlier.
\end{chapteropening}

\section[chā]{\bn{চা} — tea, read}

\label{lesson:BN-W04-chaa-read}



\begin{quote}
Write \textbf{\bn{চ}}, and say where in the mouth it is made.
\end{quote}

\subsection*{Script}
\begin{quote}
\textbf{\bn{চ}} + \textbf{\bn{া}} $\to$ \textbf{\bn{চা}}
\end{quote}

One consonant and one sign. \textbf{chā} — \emph{tea}, and the shortest word this book will ask anyone to read.

It is also a loanword, and the shape of the loan is visible in the sound. Chinese \emph{chá} travelled overland through Persian \emph{čā} and arrived in Bengali as \textbf{\bn{চা}}. The same leaf travelled by sea, through Hokkien and Dutch traders, and arrived in English as \emph{tea}. One plant, two roads, two words — and which word a language uses tells you which road the leaf took to get there.

Two pieces, and both of them arrived long before this page: \textbf{\bn{চ}} from the consonant chapter, \textbf{\bn{া}} from the very first script chapter. Nothing here is new except the permission to put them together — which is what a reading lesson is for.

\subsection*{Your turn}
\begin{itemize}
\raggedright
  \item \emph{Read it:} \bn{চা}
  \item \emph{Write it:} \bn{চা}
\end{itemize}

\begin{tcolorbox}[breakable,colback=teal!4,colframe=teal!35!black,title={Before you move on}]
How many pieces are in \textbf{\bn{চা}}? (\textbf{Two}.) And why does English say \emph{tea} where Bengali says \emph{chā}? (The word came \textbf{by sea}, not overland.)

Source: \href{https://www.unicode.org/charts/PDF/U0980.pdf}{Unicode Bengali chart}.
\end{tcolorbox}
\section[jôl]{\bn{জল} — water, read}

\label{lesson:BN-W04-jol-read}



\begin{quote}
Write \textbf{\bn{জ}}, and say what separates it from \textbf{\bn{চ}}.
\end{quote}

\subsection*{Script}
\begin{quote}
\textbf{\bn{জ}} + \textbf{\bn{ল}} $\to$ \textbf{\bn{জল}}
\end{quote}

Two letters. \textbf{No signs at all.}

\textbf{jôl} — \emph{water}. And the vowel in the middle of it, the \textbf{ô} you can plainly hear, is written nowhere. It is on the page because \textbf{\bn{জ}} carries it, the way every bare Bengali consonant does.

This is the shortest possible demonstration of the word \textbf{abugida}, which the very first script lesson used and this word finally earns. An alphabet writes every vowel. An abugida writes a consonant that already has one, and marks only the vowels that are different. \textbf{\bn{জল}} is what that looks like when nothing is different.

Set it beside the drink from earlier in this chapter:

\begin{quote}
\textbf{\bn{চা}} — one consonant, one sign \textbf{\bn{জল}} — two consonants, no signs
\end{quote}

Both are two syllables in the mouth and neither spells them the way an English reader would expect. The chapter before this one noted that Bengali \emph{eats} its drinks — \emph{jôl khāwā}, \textquotedblleft{}to eat water\textquotedblright{} — so these two words keep company in more than one way.

\subsection*{Your turn}
\begin{itemize}
\raggedright
  \item \emph{Read it:} \bn{জল}
  \item \emph{Read it:} \bn{চা} \bn{জল}
  \item \emph{Write it:} \bn{জল}
\end{itemize}

\begin{tcolorbox}[breakable,colback=teal!4,colframe=teal!35!black,title={Before you move on}]
How many vowel signs are in \textbf{\bn{জল}}? (\textbf{None}.) So where does the \textbf{ô} come from? (\textbf{The consonant carries it}.) And what does the word mean? (\emph{\textbf{Water}}.)

Source: \href{https://www.unicode.org/charts/PDF/U0980.pdf}{Unicode Bengali chart}.
\end{tcolorbox}
\section[dudh]{\bn{দুধ} — milk, and the three words this chapter gave back}

\label{lesson:BN-W03-dudh-read}



\begin{quote}
Write \textbf{\bn{দ}} and \textbf{\bn{ধ}}, and say what separates them.
\end{quote}

\subsection*{Script}
\begin{quote}
\textbf{\bn{দ}} + \textbf{\bn{ু}} + \textbf{\bn{ধ}} $\to$ \textbf{\bn{দুধ}}
\end{quote}

\begin{itemize}
\raggedright
  \item \textbf{\bn{দ}} with \textbf{\bn{ু}} under it — \emph{du}, voice on, no breath
  \item \textbf{\bn{ধ}} — voice on, \textbf{with} breath, and its own vowel gone at the end of the word
\end{itemize}

\textbf{dudh} — \emph{milk}. Two letters that a reader half-attending will read as the same shape twice, and the word says them differently: the first closes clean, the second opens onto air. Say it slowly enough to hear the second one leak.

This is the whole point of a script that spells the breath. English writes \emph{d} for both and lets speakers sort it out; Bengali writes the difference down, which is more to learn and less to guess.

Count what this chapter added. \textbf{No new pieces at all} — and three more words handed to the eye:

\begin{itemize}
\raggedright
  \item \textbf{\bn{চা}} — \emph{tea}
  \item \textbf{\bn{জল}} — \emph{water}
  \item \textbf{\bn{দুধ}} — \emph{milk}
\end{itemize}

Twenty-five pieces, ten words. That is the point of a reading chapter: the eight consonants of the last script chapter were bought once and are still paying, and the three words on this page cost nothing new to read. Two of the three turn on the breathy pairs that chapter drilled, which is the distinction an English-speaking reader has no habit for at all.

\subsection*{Your turn}
\begin{itemize}
\raggedright
  \item \emph{Read it:} \bn{চা} \bn{জল} \bn{দুধ}
  \item \emph{Write it:} \bn{দুধ}
  \item \emph{Write it:} \bn{চা} \bn{জল} \bn{দুধ}
\end{itemize}

\begin{tcolorbox}[breakable,colback=teal!4,colframe=teal!35!black,title={Before you move on}]
Which letter of \textbf{\bn{দুধ}} lets the breath out? (\textbf{The second}.) How many pieces did this chapter add? (\textbf{None}.) And how many words can you read now? (\textbf{Ten}.)

Source: \href{https://www.unicode.org/charts/PDF/U0980.pdf}{Unicode Bengali chart}.
\end{tcolorbox}
